\documentclass[11pt]{article}
\usepackage[a4paper,margin=28mm]{geometry}
\usepackage{amsmath,amssymb,amsthm}
\usepackage[T1]{fontenc}
\usepackage{lmodern}
\usepackage{microtype}
\usepackage[colorlinks=true,linkcolor=blue,citecolor=blue,urlcolor=blue]{hyperref}
\newtheorem{theorem}{Theorem}
\newtheorem{lemma}[theorem]{Lemma}
\newtheorem{corollary}[theorem]{Corollary}
\newcommand{\Pf}{\mathcal P_{\mathrm{fin}}}
\newcommand{\HH}{\mathcal H}
\title{Counterexamples to Wang's Conjectures 3.26 and 5.3\protect\\ on semigroup algebras}
\author{Haobo Ma\thanks{ChronoAI Pte Ltd, Singapore.}
\and Wenlin Zhang\thanks{National University of Singapore; ChronoAI Pte Ltd, Singapore.}}
\date{1 November 2026}
\begin{document}
\maketitle
\begin{abstract}
We refute Conjectures 3.26 and 5.3 of Wang's
\emph{Semigroup rings and algebraically independent sequences with respect
to idempotents in commutative semigroups}
\cite[arXiv v1, pp.~23 and 27]{Wang2025}.
The commutative periodic monoid
$\Pf(\mathbb N)\times\{0,1,2\}$, with union in the first coordinate and
addition capped at two in the second, has algebraic Davenport constant one
over every nonzero commutative unitary ring. Its nongroup relative Davenport
constants are unbounded. The finite submonoids with supports contained in
$\{1,\ldots,m\}$, $m\geq1$, have the same algebraic constant,
whereas the expression conjectured to equal it grows without bound as $m$ increases.
\end{abstract}

\section{The question and the construction}

Wang's Conjecture 3.26 \cite[p.~23]{Wang2025} asserts that finiteness of
$d(S,R)$ forces bounded relative Davenport constants in the Green quotient
at elements whose positive cyclic semigroups are not groups. We construct
a counterexample with $d(S,R)=1$. The same construction gives finite
counterexamples to the equality in Conjecture 5.3 \cite[p.~27]{Wang2025}.
The mechanism is that differences between transient elements and their
cyclic idempotents multiply to zero, while unions retain arbitrarily many
indispensable support elements.

The declaration \texttt{result} in the cited \emph{trureturing} source
\cite{Trureturing} negates Conjecture~3.26 over $\mathbb Q$, using
$d\leq1$ and unbounded nongroup relative constants. The proofs below
establish the arbitrary-ring equality, all relative constants, and the
finite counterexamples to Conjecture~5.3.

We use Wang's definitions of $d(S,R)$ \cite[Definition~2.1, p.~8]{Wang2025}
and $D_a(S)$ \cite[Definition~D, p.~9]{Wang2025}; the latter invariant is
studied in \cite{Wang2017}. Even for finite groups, $d(G,K)$ can exceed
$D(G)-1$ \cite[Theorem~1.1]{Smertnig2010}.
Here the discrepancy occurs with trivial Sch\"utzenberger groups.
General treatments of the underlying semigroups and semigroup rings are
\cite{Grillet2001,Gilmer1984}.

Let $S$ be a commutative periodic semigroup, written additively, and let
$R$ be a nonzero commutative ring with identity. The cyclic semigroup
$\langle s\rangle=\{ns:n\geq1\}$ has a unique idempotent $e(s)$.
The \emph{ordinary} semigroup algebra $R[S]$ is free over $R$ on all basis
elements $X^s$, with $X^sX^t=X^{s+t}$. In particular, an absorbing
semigroup element, if present, remains a basis label.
For a finite sequence $T=(s_1,\ldots,s_\ell)$, set
\[
 F_T(c_1,\ldots,c_\ell)
 =\prod_{i=1}^{\ell}(X^{s_i}-c_iX^{e(s_i)}).
\]
Then $d(S,R)$ is the supremum of the lengths of sequences for which
$F_T(c_1,\ldots,c_\ell)\ne0$ for \emph{every} choice of nonzero
$c_1,\ldots,c_\ell\in R$.

The relative constant $D_a(S)$ is the supremum of lengths of sequences
with sum $a$ having no proper subsequence with the same sum. In a monoid
the empty sequence has sum $0_S$ and is allowed as a subsequence.
Write $(a)=\{a\}\cup(a+S)$ for the principal ideal and
$a\HH b$ when $(a)=(b)$. Write $\bar a$ for the image of $a$ in
$S/\HH$.

For a set $U$, put
\begin{equation}\label{eq:model}
 S_U=\Pf(U)\times\{0,1,2\},\qquad
 (A,i)+(B,j)=(A\cup B,\min\{2,i+j\}).
\end{equation}
Its identity is $(\varnothing,0)$.

\begin{theorem}\label{thm:main}
For every set $U$ and every nonzero commutative unitary ring $R$,
$S_U$ is a commutative periodic monoid and
\begin{equation}\label{eq:exact}
 d(S_U,R)=1,\qquad \HH=\Delta_{S_U},\qquad
 D_{(A,k)}(S_U)=|A|+k\quad(0\leq k\leq2).
\end{equation}
The positive cyclic semigroup of $(A,k)$ is not a group exactly when
$k=1$. Every Sch\"utzenberger group of $S_U$ is trivial.
Consequently $S_{\mathbb N}$ refutes Conjecture~3.26.
\end{theorem}

\section{Proof of the exact formulas}

Associativity of \eqref{eq:model} follows from associativity of union and
\[
 \min\{2,\min\{2,i+j\}+k\}=\min\{2,i+j+k\}.
\]
Commutativity and the identity follow directly. For $x=(A,k)$ and $n\geq1$,
\[
 nx=(A,\min\{2,nk\}).
\]
Thus levels zero and two are idempotent. At level one the positive cyclic
semigroup consists of $x=(A,1)$ and $e(x)=(A,2)$, with
$x+x=x+e(x)=e(x)+e(x)=e(x)$. It has no identity: neither of its two
elements fixes $x$. Singleton cyclic semigroups at the other levels are
groups with their own identities. This proves periodicity and the
nongroup classification.

\begin{lemma}\label{lem:algebra}
For every $A,B\in\Pf(U)$,
\[
 (X^{(A,1)}-X^{(A,2)})(X^{(B,1)}-X^{(B,2)})=0.
\]
Moreover $d(S_U,R)=1$.
\end{lemma}
\begin{proof}
All four monomials in the displayed product have label $(A\cup B,2)$,
and their coefficients sum to $1-1-1+1=0$.
Consider any sequence of length at least two and choose every coefficient
$c_i=1_R$. If it contains an idempotent, its corresponding factor is zero.
Otherwise all its elements have level one and the first two factors
multiply to zero. Hence no such sequence satisfies the universal
nonvanishing requirement, proving $d(S_U,R)\leq1$.
For $x=(\varnothing,1)$ and any nonzero $c\in R$,
$X^x-cX^{e(x)}\ne0$: the two labels differ and the coefficient at $x$
is $1_R\ne0$. This proves the reverse inequality.
\end{proof}

\begin{lemma}\label{lem:green}
For $a=(A,k)\in S_U$,
\[
 (a)=\{(B,j): A\subseteq B\in\Pf(U),\ k\leq j\leq2\}.
\]
Consequently $\HH$ is equality and every Sch\"utzenberger group is trivial.
\end{lemma}
\begin{proof}
Addition can only enlarge the support and increase the level. Conversely,
if $A\subseteq B$ and $k\leq j\leq2$, then
$(A,k)+(B,j-k)=(B,j)$. Equality of two principal ideals therefore forces
both inclusions of their supports and both inequalities of their levels,
so their generators coincide. Each Green class is a singleton; its
stabilizer is nonempty because it contains $0_{S_U}$, and all its induced
permutations are the identity. Its Sch\"utzenberger group is thus the
trivial group.
\end{proof}

\begin{lemma}\label{lem:relative}
For every $A\in\Pf(U)$ and $k\in\{0,1,2\}$,
$D_{(A,k)}(S_U)=|A|+k$.
\end{lemma}
\begin{proof}
Take $k$ copies of $(\varnothing,1)$ together with one term
$(\{a\},0)$ for each $a\in A$. The sum is $(A,k)$.
Deleting any singleton-support term loses that support element; deleting
any of the $k$ marked terms lowers the level. Every proper subsequence
therefore changes the sum. This also covers $A=\varnothing,k=0$:
the empty sequence is irreducible of length zero.

For the upper bound let $T$ be irreducible with sum $(A,k)$.
If $k=0$, mark no terms. If $k=1$, there is exactly one term of level one
and no term of level two; mark that term. If $k=2$, mark one level-two
term when one occurs, and otherwise mark two level-one terms.
In each case the marked terms alone supply level $k$, and their number
is at most $k$. For every unmarked term $(B,j)$, deleting it preserves
the level. Irreducibility forces $B$ to contain a point absent from
the supports of all other terms. Such exclusive points for different
unmarked terms are distinct elements of $A$. There are therefore at
most $|A|$ unmarked terms, giving $|T|\leq |A|+k$.
\end{proof}

The lemmas prove \eqref{eq:exact}. In $S_{\mathbb N}$ the targets
$a_m=(\{1,\ldots,m\},1)$ have nongroup positive cyclic semigroups and
\[
 D_{\bar a_m}(S_{\mathbb N}/\HH)=m+1.
\]
This is an unbounded sequence of finite relative constants in one fixed
infinite periodic monoid. Together with $d(S_{\mathbb N},R)=1$, it
contradicts Conjecture~3.26.

\section{The finite conjecture and the companion remark}

For $a\in S$ let $\epsilon_a=1$ if $\langle a\rangle$ is not a group,
and let $\epsilon_a=0$ otherwise. Wang's Conjecture~5.3 concerns finite
commutative $S$ and a field $K$ that is either a splitting field for $S$
or has characteristic $p$ with $\exp(S)$ a $p$-power. It proposes
\begin{equation}\label{eq:conjectured}
 d(S,K)=\max_{a\in S}
 \{\epsilon_aD_{\bar a}(S/\HH)+d(\Gamma(H_a),K)\},
\end{equation}
where $\Gamma(H_a)$ is the Sch\"utzenberger group.

\begin{corollary}\label{cor:finite}
Let $m\geq1$ and $S_m=S_{\{1,\ldots,m\}}$. For every field $K$,
\[
 |S_m|=3\cdot2^m,\qquad \exp(S_m)=1,\qquad d(S_m,K)=1,
\]
whereas the right side of \eqref{eq:conjectured} is $m+1$.
Thus Conjecture~5.3 is false, already for a six-element monoid.
\end{corollary}
\begin{proof}
Every cyclic period is one, so $\exp(S_m)=1$. By Wang's definition
\cite[p.~9]{Wang2025}, a splitting field for $S$ has exactly $\exp(S)$
roots of $z^{\exp(S)}=1$. Every field satisfies this condition when
$\exp(S)=1$; thus the first hypothesis of the conjecture applies.
Theorem~\ref{thm:main} gives $d(S_m,K)=1$ and trivial
$\Gamma(H_a)$. For the trivial group the algebraic constant is zero,
since coefficient one makes any nonempty sequence's factor zero.
Only level-one elements have $\epsilon_a=1$, and their relative
constants are $|A|+1$. Their maximum is $m+1$, attained at
$(\{1,\ldots,m\},1)$. For $m=1$ this gives $1\ne2$.
\end{proof}

The same formulas also give $D(S_m)=m+3$, with $D(S)$ defined as the
least length from which every sequence is reducible: the largest
irreducible length is $m+2$, attained at $(\{1,\ldots,m\},2)$.
Thus the gap in \eqref{eq:conjectured} is unbounded even in this finite
family with constant algebraic Davenport invariant.

Remark~3.27 of \cite[p.~23]{Wang2025} conditionally proposes a finiteness
characterization after Conjecture~3.26. Its printed wording says
``group'', rather than ``semigroup''. For the semigroup interpretation,
our infinite example disproves the necessity of the bounded-relative-constant
condition, even though every Sch\"utzenberger group has algebraic
constant zero. The literal group statement has no nongroup positive
cyclic elements and is not contradicted by this example.

\begin{thebibliography}{6}
\bibitem{Wang2025}
G.~Wang, \emph{Semigroup rings and algebraically independent sequences
with respect to idempotents in commutative semigroups},
\href{https://arxiv.org/abs/2509.18724v1}{arXiv:2509.18724v1}, 2025.

\bibitem{Wang2017}
G.~Wang, \emph{Additively irreducible sequences in commutative semigroups},
J. Combin. Theory Ser. A \textbf{152} (2017), 380--397.
\href{https://doi.org/10.1016/j.jcta.2017.07.001}{doi:10.1016/j.jcta.2017.07.001}.

\bibitem{Smertnig2010}
D.~Smertnig, \emph{On the Davenport constant and group algebras},
Colloq. Math. \textbf{121} (2010), no.~2, 179--193.
\href{https://doi.org/10.4064/cm121-2-2}{doi:10.4064/cm121-2-2}.

\bibitem{Grillet2001}
P.~A.~Grillet, \emph{Commutative Semigroups}, Advances in Mathematics,
vol.~2, Kluwer Academic Publishers, Dordrecht, 2001.

\bibitem{Gilmer1984}
R.~Gilmer, \emph{Commutative Semigroup Rings}, Chicago Lectures in
Mathematics, University of Chicago Press, Chicago, 1984.

\bibitem{Trureturing}
The trureturing contributors, \emph{Formal source for Wang's Conjecture 3.26},
\href{https://github.com/the-omega-institute/trureturing/blob/d1876cfbbef8c13023311126691b180eefa69569/D5/S0/Certificates/Algebra/WangConjecture326Refutation.lean}{\texttt{WangConjecture326Refutation.lean}},
\emph{trureturing}, source revision of 31 October 2026.
\end{thebibliography}

\smallskip
\noindent\small This work is licensed under the Creative Commons
Attribution 4.0 International License (CC BY 4.0).
\end{document}
